\setcounter{section}{9}

  \zwischenueberschrift{Soal pemanasan}

\inputaufgabe
{}
{

Tunjukkan bahwa setiap
\definitionsverweis {vektor singgung}{}{}

\mavergleichskette
{\vergleichskette
{ v
}
{ \in }{ T_PS^2
}
{  }{
}
{    }{
}
{   }{
}
}
{}{}{}
pada suatu titik $P$ di
\definitionsverweis {permukaan bola satuan}{}{}
dapat direalisasikan oleh sebuah lintasan
\anfuehrung{berkecepatan tetap}{}
yang
\definitionsverweis {diferensiabel}{}{}
pada sebuah
\definitionsverweis {lingkaran besar}{}{.}

}
{} {}

\inputaufgabe
{}
{

Berikan realisasi sesederhana mungkin bagi
\definitionsverweis {vektor-vektor singgung}{}{}
pada suatu titik
\mathl{P=(s,t)}{} di permukaan silinder
\mathl{S^1 \times \R}{.}

}
{} {}

\inputaufgabe
{}
{

Misalkan
\mathkor {} {M} {dan} {N} {}
adalah
\definitionsverweis {manifold diferensiabel}{}{}
dan
\maabbdisp {\varphi} { M } { N
} {}
sebuah
\definitionsverweis {pemetaan diferensiabel}{}{.}
Misalkan $P \in M$ dan
\mathl{Q=\varphi(P)}{}, serta misalkan
\maabbdisp {\gamma_1,\gamma_2} { I } { M
} {}
dua
\definitionsverweis {kurva terdiferensialkan}{}{}
dengan interval terbuka
\mathl{0 \in I}{} dan
\mathl{\gamma_1 (0) = \gamma_2(0)=P}{.} Misalkan
\mathkor {} {\gamma_1} {dan} {\gamma_2} {}
\definitionsverweis {ekuivalen secara tangensial}{}{}
di titik $P$.
Tunjukkan bahwa
\definitionsverweis {komposisi}{}{}
\mathkor {} {\varphi \circ \gamma_1} {dan} {\varphi \circ \gamma_2} {}
juga ekuivalen secara tangensial di $Q$.

}
{} {}

\inputaufgabe
{}
{

Berikan sebuah contoh pemetaan yang
\definitionsverweis {injektif}{}{,}
tidak
\definitionsverweis {surjektif}{}{,}
dan
\definitionsverweis {diferensiabel}{}{},
\maabbdisp {\varphi} {\R} {\R
} {}
sedemikian sehingga
\definitionsverweis {pemetaan singgung}{}{}
yang bersesuaian,
\maabbdisp {T_P(\varphi)} {T_P \R} {T_{\varphi(P)} \R
} {}
bijektif pada setiap titik

\mavergleichskette
{\vergleichskette
{ P
}
{ \in }{ \R
}
{  }{
}
{    }{
}
{   }{
}
}
{}{}{.}

}
{} {}

\inputaufgabe
{}
{

Berikan sebuah contoh pemetaan yang
\definitionsverweis {surjektif}{}{,}
tidak
\definitionsverweis {injektif}{}{,}
dan
\definitionsverweis {diferensiabel}{}{},
\maabbdisp {\varphi} {S^1} {S^1
} {}
sedemikian sehingga
\definitionsverweis {pemetaan singgung}{}{}
yang bersesuaian,
\maabbdisp {T_P(\varphi)} {T_P S^1} {T_{\varphi(P)} S^1
} {}
bijektif pada setiap titik
\mathl{P \in S^1}{.}

}
{} {}

\inputaufgabe
{}
{

Tunjukkan bahwa pada sebuah
\definitionsverweis {manifold diferensiabel}{}{}, penjumlahan lintasan
\maabbdisp {\gamma_1, \gamma_2} { I } { M
} {}
dengan
\mathl{\gamma_1(0) =\gamma_2(0)=P \in M}{,} yang dapat diperoleh melalui sebuah peta
\maabbdisp {\alpha} { U } { V
} {}
dengan
\mathl{\alpha(P)=0}{} dari penjumlahan di $\R^n$, pada umumnya bergantung pada peta yang dipilih.

}
{} {}

\inputaufgabe
{}
{

Tunjukkan bahwa pemetaan singgung $T_P(\varphi)$ yang bersesuaian dengan
\maabbeledisp {\varphi} {\R^1} {S^1
} { t } {( \cos  t, \sin  t  )
} {,}
bijektif pada setiap titik
\mathl{P \in \R}{.}

}
{} {}

\inputaufgabe
{}
{

Berikan sebuah contoh pemetaan yang
\definitionsverweis {surjektif}{}{}
dan
\definitionsverweis {diferensiabel}{}{}
\maabbdisp {\varphi} { M } { N
} {}
antara dua
\definitionsverweis {manifold diferensiabel}{}{},
\mathkor {} {M} {dan} {N} {}
sedemikian sehingga
\definitionsverweis {pemetaan singgung}{}{}
yang bersesuaian,
\maabbdisp {T_P(\varphi)} {T_PM} {T_{\varphi(P)}N
} {}
tidak surjektif pada suatu titik
\mathl{P \in M}{.}

}
{} {}

\inputaufgabe
{}
{

Berikan sebuah contoh pemetaan yang
\definitionsverweis {injektif}{}{}
dan
\definitionsverweis {diferensiabel}{}{}
\maabbdisp {\varphi} { M } { N
} {}
antara dua
\definitionsverweis {manifold diferensiabel}{}{},
\mathkor {} {M} {dan} {N} {}
sedemikian sehingga
\definitionsverweis {pemetaan singgung}{}{}
yang bersesuaian,
\maabbdisp {T_P(\varphi)} {T_PM} {T_{\varphi(P)}N
} {}
tidak injektif pada suatu titik
\mathl{P \in M}{.}

}
{} {}

\inputaufgabe
{}
{

Misalkan $M$ sebuah
\definitionsverweis {manifold diferensiabel}{}{}
dan

\mavergleichskette
{\vergleichskette
{ P
}
{ \in }{ M
}
{  }{
}
{    }{
}
{   }{
}
}
{}{}{}
sebuah titik. Perhatikan himpunan berikut.

\mavergleichskettedisp
{\vergleichskette
{ T
}
{ =} { \left\{  (U,f)   \mid   U \subseteq M \text{ terbuka} , \,  P \in U   , \,  f \in C^1(U,\R)    \right\}
}
{ } {
}
{ } {
}
{ } {
}
}
{}{}{.}
Perhatikan
\definitionsverweis {relasi}{}{}

\mathdisp {(U,f) \sim (V,g) : \text{ terdapat himpunan terbuka } W \text{ dengan } P \in W \subseteq U \cap V \text{ dan } f  \vert_W = g \vert_W} { . }

\aufzaehlungzwei {Tunjukkan bahwa relasi ini merupakan sebuah
\definitionsverweis {relasi ekuivalensi}{}{}
pada $T$.
} {Tunjukkan bahwa terdapat suatu
\definitionsverweis {struktur gelanggang}{}{}
alami pada himpunan
\definitionsverweis {kelas-kelas ekuivalensi}{}{}
dari relasi ekuivalensi tersebut.
}

}
{} {}

\inputaufgabe
{}
{

Misalkan $X$ sebuah
\definitionsverweis {ruang topologis}{}{,}
$Y$ sebuah himpunan, dan
\maabbdisp {\varphi} { X } { Y
} {}
sebuah
\definitionsverweis {pemetaan}{}{.}
Tunjukkan bahwa
\definitionsverweis {sistem himpunan}{}{}

\mathdisp {{\mathcal T } = \left\{ V \subseteq Y  \mid  \varphi^{-1}(V) \text{ terbuka di } X        \right\}} {  }

mendefinisikan sebuah
\definitionsverweis {topologi}{}{}
pada $Y$ sehingga $\varphi$
\definitionsverweis {kontinu}{}{.}

}
{} {}

Topologi yang diperkenalkan dalam soal sebelumnya disebut \stichwort {topologi citra} {.}

\inputaufgabe
{}
{

Tunjukkan bahwa pada $\R^n$, hubungan

\mathdisp {P \sim Q, \text{ jika } P-Q \in \Z^n} {  }

mendefinisikan sebuah
\definitionsverweis {relasi ekuivalensi}{}{.}
Lengkapi
\definitionsverweis {himpunan hasil bagi}{}{}

\mavergleichskettedisp
{\vergleichskette
{Y
}
{ =} {\R^n /\!\!\sim
}
{ =} {\R^n/\Z^n
}
{ } {
}
{ } {
}
}
{}{}{}
dengan
\definitionsverweis {topologi citra}{}{}
yang diinduksi oleh
\definitionsverweis {pemetaan hasil bagi}{}{}
\maabb {\varphi} {\R^n } { Y
} {}
tersebut. Tunjukkan bahwa $Y$ merupakan sebuah
\definitionsverweis {ruang Hausdorff}{}{.}

}
{} {}

\inputaufgabe
{}
{

Uraikan garis besar suatu teori tentang \anfuehrung{manifold kompleks}{.} Apa manifold riil yang mendasarinya, berapakah
\zusatzklammer {kompleks/riil} {} {}
dimensinya, dan seperti apakah ruang singgungnya?

}
{} {}

  \zwischenueberschrift{Soal untuk dikumpulkan}

\inputaufgabe
{4}
{

Misalkan $M$ sebuah
\definitionsverweis {manifold diferensiabel}{}{} dan
\mathl{P \in M}{.} Tunjukkan bahwa untuk sebuah
\definitionsverweis {kurva terdiferensialkan}{}{}
\maabbdisp {\gamma} { I } { M
} {}
dengan
\mathl{\gamma(0)= P}{} dan
\mathl{a \in \R}{}, di dalam
\definitionsverweis {ruang singgung}{}{}
\mathl{T_PM}{} berlaku hubungan

\mavergleichskettedisp
{\vergleichskette
{a [\gamma]
}
{ =} { [  \lambda]
}
{ } {
}
{ } {
}
{ } {
}
}
{}{}{}. Definisikan interval terbuka
\mathl{J:=\{t\in\R\mid at\in I\}}{}
dan pemetaan $\lambda:J\to M$ oleh
\mathl{\lambda(t):= \gamma(at)}{}; tunjukkan hubungan yang sama untuk kelas kurva ini.

Catatan edisi: Sumber mendefinisikan $\lambda$ pada interval semula $I$, padahal $at$ tidak harus berada di $I$. Domain $J$ di atas adalah prapeta $I$ oleh perkalian dengan $a$, selalu merupakan lingkungan terbuka dari nol, dan membuat reparametrisasi terdefinisi dengan benar.

}
{} {}

\inputaufgabe
{6}
{

Misalkan $M$ sebuah
$C^k$-\definitionsverweis {manifold}{}{}
dan

\mavergleichskette
{\vergleichskette
{ P
}
{ \in }{ M
}
{  }{
}
{    }{
}
{   }{
}
}
{}{}{.}
Untuk
$C^k$-\definitionsverweis {kurva}{}{}
\maabbdisp {\gamma_1,\gamma_2} { I } { M
} {}
dengan

\mavergleichskette
{\vergleichskette
{ \gamma_1(0)
}
{ = }{ \gamma_2(0)
}
{ = }{ P
}
{    }{
}
{   }{
}
}
{}{}{}, definisikan sebuah relasi ekuivalensi yang, dalam
\zusatzklammer {setiap} {} {}
\definitionsverweis {peta}{}{,}
memperhitungkan
\definitionsverweis {turunan}{}{}
hingga orde $k$. Seperti apakah wakil-wakil sederhana dari relasi ekuivalensi ini? Untuk $k=1$, pulihkan struktur ruang vektor pada
\definitionsverweis {himpunan hasil bagi}{}{}
dan tentukan dimensinya. Untuk $k\geq 2$, periksa apakah konstruksi ruang vektor itu tidak bergantung pada peta. Setelah sebuah peta dipilih, identifikasikan himpunan hasil bagi dengan $(\R^n)^k$ dan tentukan dimensinya sebagai manifold; lalu tunjukkan melalui suatu pergantian peta nonlinier bahwa struktur ruang vektor tersebut pada umumnya tidak kanonis.

Catatan edisi: Perintah terakhir dalam sumber meminta struktur ruang vektor untuk semua $k$. Untuk $k\geq2$, turunan berorde lebih tinggi berubah secara nonlinier di bawah pergantian peta, sehingga ruang jet kurva ini berdimensi $kn$ sebagai manifold tetapi tidak mempunyai struktur ruang vektor kanonis tanpa pilihan tambahan seperti peta atau koneksi.

}
{} {}

\inputaufgabe
{5}
{

Misalkan $M$ sebuah
\definitionsverweis {manifold diferensiabel}{}{}
dan
\mathl{P \in M}{.} Kita mengatakan bahwa dua
\definitionsverweis {kurva}{}{}
\maabbdisp {\gamma_1, \gamma_2} { I } { M
} {}
dengan
\mathl{\gamma_1(0)= \gamma_2(0)= P}{} mendefinisikan \stichwort {germ kurva} {} yang sama jika terdapat
\mathl{\epsilon >0}{} sedemikian sehingga

\mavergleichskettedisp
{\vergleichskette
{ \gamma_1 \!\mid_{ [-\epsilon, \epsilon]}
}
{ =} { \gamma_2 \!\mid_{ [-\epsilon, \epsilon]}
}
{ } {
}
{ } {
}
{ } {
}
}
{}{}{.}

a) Tunjukkan bahwa ini mendefinisikan sebuah
\definitionsverweis {relasi ekuivalensi}{}{}
pada himpunan semua kurva
\maabb {\gamma} { I } { M
} {}
dengan $\gamma(0)= P$
\zusatzklammer {dan dengan interval terbuka berbeda-beda
\mathl{0 \in I}{}} {} {.}

b) Tunjukkan bahwa
\definitionsverweis {kurva-kurva terdiferensialkan}{}{,}
yang merepresentasikan germ kurva yang sama, juga merepresentasikan vektor singgung yang sama.

}
{} {}

\inputaufgabe
{8}
{

Misalkan
\definitionsverweis {ruang hasil bagi}{}{}

\mavergleichskette
{\vergleichskette
{ Y
}
{ = }{ \R^n/\Z^n
}
{  }{
}
{    }{
}
{   }{
}
}
{}{}{}
dilengkapi dengan
\definitionsverweis {topologi citra}{}{.}
Definisikan sebuah struktur manifold pada $Y$ dengan peta-peta yang sesuai. Tunjukkan bahwa
\definitionsverweis {pemetaan hasil bagi}{}{}
\maabbdisp {\varphi} {\R^n } { Y
} {}
merupakan
\definitionsverweis {pemetaan diferensiabel}{}{,}
dan bahwa
\definitionsverweis {pemetaan singgung}{}{}
merupakan
\definitionsverweis {isomorfisme}{}{}
di setiap titik.

}
{} {}

 [[Kategorie:Latexseite]]
